\footnotesize
\arrayrulecolor{black!22}
\begin{longtable}{L{0.04\linewidth} L{0.105\linewidth} L{0.135\linewidth} L{0.145\linewidth} L{0.12\linewidth} L{0.115\linewidth} L{0.115\linewidth} L{0.12\linewidth} L{0.075\linewidth}}
\hline
\textbf{ID} & \textbf{Goal} & \textbf{Motivation} & \textbf{Vision} & \textbf{Raw data /\allowbreak{} tools} & \textbf{Operational scope} & \textbf{Not-before gate} & \textbf{Derived-only boundary} & \textbf{Plan action}\\
\hline
\endfirsthead
\hline
\textbf{ID} & \textbf{Goal} & \textbf{Motivation} & \textbf{Vision} & \textbf{Raw data /\allowbreak{} tools} & \textbf{Operational scope} & \textbf{Not-before gate} & \textbf{Derived-only boundary} & \textbf{Plan action}\\
\hline
\endhead
SLTG-01 & Semi-automated documentation of existing Lean code inside the same planning/\allowbreak{}proofbook tool. & The Bestandscode cannot be fully normalized at once, but every module, declaration family, object, proof bridge, and consumer path can be documented with the same YAML/\allowbreak{}proof-dossier infrastructure so later interpretation is minimized. & Existing Lean-code documentation should grow into a reviewable object universe: each missing object/\allowbreak{}quantity becomes a dossier-backed node that can be inspected by the relevant expert without reading the full repository. & repo\_inventory/\allowbreak{}raw\_export modules.json, import\_edges.json, symbol\_reference\_summary.json, dependency\_order, mainline forensics, object register, proof dossiers, implementation scratchpad. & Secondary workflow: document existing code and its true consumers while Phase 1.x cleans the public API. It does not authorize refactors or deletion by itself. & After Phase 1.1.1.1.4 declaration consumption classification and initial proof-dossier infrastructure are closed;\newline deeper execution after Phase 1.1.5 import firewall is green. & Documentation may classify provenance and consumers;\newline it may not turn graph metadata into CNNA-generative assumptions. & Add Phase 19.1 and object LEG16;\newline use the generated Lean module graph to prevent forgotten code/\allowbreak{}documentation gaps. V1.12 adds LEG36 so existing-code documentation also records semantic consolidation/\allowbreak{}no-merge candidates before new substantive code.\\
\hline
\addlinespace[0.24em]
SLTG-02 & Derivation-graph geometry of the formalized theory. & The repository already contains graph-theoretic data: import graph, dominators, layer decomposition, source-sink signal paths, and robust semantic mainline. These can later expose bottlenecks, weak derivation edges, and discipline clusters. & A later graph/\allowbreak{}database layer can provide semantic zoom from pillars to phases, modules, declarations, objects, proof dossiers, and theorem sites;\newline the graph remains a diagnostic/\allowbreak{}evidence map, not a generator. & import\_edges.json, dominator analysis, layer decomposition, source-to-sink paths, module group graph, architectural core summary, mainline forensics. & Secondary research/\allowbreak{}diagnostic layer only. Use it to understand theory organization and proof pressure after core derivations are stable. & After the neutral API, generator governance, object proof dossiers, and main operational arithmetic/\allowbreak{}DtN paths have stabilized;\newline not before active Phase 1.x cleanup. & Graph geometry can suggest audits and consistency checks but cannot supply mathematical axioms, stochasticity, open-system behavior, or physical interpretation. & Add Phase 19.2 and object LEG17;\newline render/\allowbreak{}import future graph metrics such as dominators, cuts, signal-path depth, cluster boundaries, and optional Cheeger-style diagnostics as comparison metadata.\\
\hline
\addlinespace[0.24em]
SLTG-03 & Formal emergence-path forensics for open-system dynamics inside the derivation network. & Later, once Pillar A/\allowbreak{}C objects are derived, the proof/\allowbreak{}module graph may help locate where closed finite subsystem assumptions become open effective dynamics through boundary/\allowbreak{}complement/\allowbreak{}regime handoffs. & When A/\allowbreak{}C open-system paths exist, graph forensics can show where regime, complement, boundary, and handoff structures make openness visible, but the dynamics must still be derived in Lean-facing CNNA objects. & A-regime profiles, complement-family profiles, DtN/\allowbreak{}Schur handoff modules, C/\allowbreak{}OQS consumer path, source-sink paths, dominator/\allowbreak{}cut analysis, object proof dossiers. & Post-core analysis of the formalization topology;\newline not a replacement for deriving open-system behavior inside CNNA. & After A regime/\allowbreak{}complement path, A->C handoff, and C OQS consumer roadmaps have enough checked objects to analyze. Treat as bonus/\allowbreak{}secondary until then. & Emergence claims must still be proved from CNNA objects;\newline graph forensics only localize proof paths, bottlenecks, and structural analogies. & Add Phase 19.3 and object LEG18;\newline keep as black/\allowbreak{}secondary target until the relevant A/\allowbreak{}C derivations exist.\\
\hline
\addlinespace[0.24em]
SLTG-04 & Reflective comparison between CNNA object-structure and formalization-structure. & If the formal derivation graph later exhibits structural features resembling CNNA objects - layering, boundary cuts, operators on graphs, signal flow, modular clusters - this should be studied as a meta-level observation, not as a premise. & The tool may be read as a meta-level interface analogue: one active cursor/\allowbreak{}bright object is examined against the large Lean-code complement;\newline this is an operational review metaphor, not a CNNA theorem. & discipline subgraphs for Dirichlet/\allowbreak{}DtN/\allowbreak{}Schur, prenumerical/\allowbreak{}ToC/\allowbreak{}Substrate, Topology/\allowbreak{}Algebra/\allowbreak{}CD;\newline module clusters;\newline graph cuts;\newline source-sink paths;\newline proof-dossier links. & Far-target meta-analysis and visualization after mathematical closure. It may produce diagnostics, figures, and research notes, not new CNNA base assumptions. & Only after core derived objects, public API, and consumer handoffs are stable enough that graph measurements are meaningful. & Strictly comparison/\allowbreak{}diagnostic. Any apparent self-similarity between theory and formalization remains an observation unless separately formalized and proved. & Add Phase 19.4 and object LEG19;\newline keep outside the active implementation path.\\
\hline
\addlinespace[0.24em]
LT5 & Thread/\allowbreak{}full-boundary derivation-geometry and string-analogy diagnostics & The IdealThread/\allowbreak{}full-cut extrema may later illuminate how one-dimensional recursive routes sit inside full boundary/\allowbreak{}operator geometry and how open-system behavior emerges from complement/\allowbreak{}interface structure. & Future string-analogy diagnostics may use derivation geometry as a comparison map around the thread/\allowbreak{}full-boundary extrema, only after the underlying route is derived and reviewed. & IdealThread/\allowbreak{}CutSpec.full profiles;\newline import graph;\newline derivation graph;\newline Schur/\allowbreak{}DtN route audit;\newline proof dossiers;\newline source-sink paths and dominators. & Secondary diagnostic/\allowbreak{}comparison study only, after the underlying CNNA objects are code-derived and public-final. & After S07, S08, BR06, BR07, S05/\allowbreak{}S06, and the relevant arithmetic/\allowbreak{}number-theory proof dossiers are closed. & The diagnostic may compare formalization geometry to CNNA objects but cannot feed back as a generative axiom, formula, or shortcut. & Keep as Phase 19.5 side-road;\newline do not block active Phase 1 cleanup or Phase 3/\allowbreak{}4 neutralization.\\
\hline
\addlinespace[0.24em]
SLTG-06 & Database-ready CNNA Explorer and review/\allowbreak{}documentation structure & A database-backed representation is likely the most effective path for querying, analyzing, visualizing, sharing, online display, interactive exploration, and discipline-separated expert review of the CNNA object world. & Long-term vision: a custom CNNA Explorer that starts as a dark, SQL/\allowbreak{}YAML-driven, visual research interface over the Single-YAML-Master mirror and can gradually grow into a Universe Explorer. The user-facing surface should support Pillar hubs, graph navigation, right-side inspectors, filters by context/\allowbreak{}theme/\allowbreak{}phase/\allowbreak{}object/\allowbreak{}module/\allowbreak{}status/\allowbreak{}edge semantics, history timelines, multilingual descriptions, hypothesis tracking and later lab views. The explorer remains generated-secondary: Lean code is the primary truth system, YAML is the curated semantic mirror, SQL is a read-only presentation mirror. & Single YAML master, generated TSV/\allowbreak{}JSON/\allowbreak{}DOT/\allowbreak{}PDF views, phase/\allowbreak{}object/\allowbreak{}proof-dossier records, module inventory, status legends, local theory context ledger, end-theory comparison ledger, future relational or graph database export.;\newline future GitHub Actions artifacts;\newline snapshot\_manifest.json;\newline Webhosting-local MySQL importer;\newline SQL dumps/\allowbreak{}local archives & Main side objective Phase 20 plus dedicated CNNA Explorer main block Phase 21. Phase 21.0 must close the completeness/\allowbreak{}coherence/\allowbreak{}refactor/\allowbreak{}source-inventory audit before Explorer schema, SQL importer or frontend implementation are treated as stable. & Principles apply immediately. Technical database export follows after YAML schema and generated views are stable enough for round-trip-safe identifiers and review links. Public deployment additionally requires backup/\allowbreak{}rollback evidence, read-only API tests and validated staging/\allowbreak{}current import. & The database is a documentation/\allowbreak{}review/\allowbreak{}query layer. It may expose, analyze, and visualize Lean/\allowbreak{}YAML-derived metadata, but it cannot supply CNNA-generative axioms, formulas, physics, stochasticity, or established-theory inputs. & P1.37 sharpens Phase 21 with 21.2.1-21.2.3: GitHub backup-gated PR automation, CI build/\allowbreak{}snapshot artifacts, and Webhosting-local MySQL staging/\allowbreak{}import/\allowbreak{}swap/\allowbreak{}rollback. The database remains read-only/\allowbreak{}generated-secondary and cannot replace Lean/\allowbreak{}YAML truth.\\
\hline
\addlinespace[0.24em]
SLTG-07 & Geometric-algebra comparison map for later CNNA-derived algebraic and spacetime structures & Quaternion/\allowbreak{}Pauli/\allowbreak{}Clifford/\allowbreak{}geometric-algebra literature provides compact target profiles for rotations, bivectors, pseudoscalars, Clifford signatures and possible spacetime/\allowbreak{}conformal comparison layers. & After CNNA derives relevant CD, boundary-algebra, AQFT/\allowbreak{}spacetime or side-comparison outputs, the tool can generate a review map showing which GA targets are matched, failed or not yet meaningful. & Horn 0709.2238v1, Hestenes references, mathlib Quaternion/\allowbreak{}CliffordAlgebra docs, Wieser/\allowbreak{}Song Lean GA formalization, proof dossiers GA0-GA6. & Long-term comparison/\allowbreak{}review map only;\newline no foundation, no phase release, no Lean generator input. & After GA0-GA6 predecessor phases expose CNNA-derived candidate data;\newline especially CD reintegration and Pillar-B boundary/\allowbreak{}local-net phases. & GA references remain blue/\allowbreak{}reference material. A target match is evidence of compatibility, not a source of truth;\newline a mismatch must be recorded rather than patched at the consumer. & V1.25 adds phases 1.4.9.*, 13.2.1, 13.6.1 and 18.5 plus objects GA0-GA6.\\
\hline
\addlinespace[0.24em]
\end{longtable}
\arrayrulecolor{black}
\normalsize
